\subsection{Post-saturation novelty refresh after claim-expansion search}

A later claim-expansion search on 2026-08-19 surfaced additional direct pressure on the interpretation of P1.D2. Preregistered Belief Revision Contracts (PBRC) already supplies an operator-agnostic protocol layer in which admissible belief change is tied to preregistered evidence triggers, allowed revision operators, priorities, fallback, and externally validated evidence witnesses \cite{pbrc2026}. Paper~I therefore does not claim generic evidence-gated belief revision or a generic control layer above revision operators. Self-Revising Discovery Systems for Science already treats scientific discovery as a verified transition between representational regimes with typed artifacts, provenance, preservation and transport obligations \cite{selfrevisingdiscovery2026}; and the Model Discovery Agent performs M-open hypothesis-class expansion after predictive inadequacy and chooses informative experiments to identify the proposed mechanism \cite{modeldiscoveryagent2026}. Generic scientific regime revision, model-class expansion, and active discrimination are therefore donor territory rather than Paper~I novelty.

The same refresh also tightens the failure-to-repair boundary. Model or Harness? formulates an interaction-centric repair-assignment problem in which the same visible failure can belong to the model, harness, environment, or grader and should be repaired at the responsible component \cite{modelorharness2026}. Diagnosis Is Not Prescription further shows that the component most causally responsible for an agent failure need not be the best intervention target \cite{diagnosisnotprescription2026}. Paper~I therefore does not claim generic component-responsibility localization, repair assignment, or the diagnosis-versus-intervention distinction. Finally, Earned Authority under a Fixed Ceiling for Evolving Agents formalizes authorization continuity across agent mutation and allows evidence-gated authority changes below an immutable effect ceiling \cite{earnedauthority2026}; generic mutation authorization and non-amplification are not Paper~I novelty either.

These sources sharpen, but do not broaden or invalidate, the frozen v2.2.4 result. They supersede any reading of the earlier phrase ``authority gate'' as a standalone generic novelty claim. The empirically tested residual is the \emph{science-specific protected necessity composition}: a high-level $K/W/M$ mutation is not licensed merely because a diagnosis names a high-level cause, because a broad intervention would locally repair the task, because an evidence trigger admits some belief change, or because a generic authorization policy permits mutation. In the frozen mechanical family, admission additionally requires that registered lower-level alternatives be excluded, the candidate restore the same task, protected siblings/invariants survive, and the observed dependency impact match the proposed reopen scope.

A separate successor study asks whether a more general revision-responsibility/minimal-escalation control layer adds value across heterogeneous domains after these donor mechanisms are granted. That study is prospective and has not produced protected outcomes. Accordingly, Paper~I makes no cross-domain revision-control superiority claim and no claim that its current mechanical result establishes a general theory of scientific self-revision.
